\documentclass[11pt,a4paper]{article}
\usepackage[margin=25mm]{geometry}
\usepackage{fontspec}
\setmainfont{TeX Gyre Pagella}
\usepackage{amsmath,amssymb,mathtools}
\usepackage{graphicx}
\usepackage{xcolor}
\usepackage[unicode,colorlinks=true,linkcolor=blue,urlcolor=blue]{hyperref}
\usepackage{bookmark}
\usepackage{enumitem}
\setlist{itemsep=3pt,topsep=5pt}
\setlength{\emergencystretch}{3em}
\newcommand{\tightlist}{\setlength{\itemsep}{0pt}\setlength{\parskip}{0pt}}
\newcommand{\passthrough}[1]{#1}
\newcommand{\pandocbounded}[1]{#1}
\makeatletter
\def\maxwidth{\ifdim\Gin@nat@width>\linewidth\linewidth\else\Gin@nat@width\fi}
\def\maxheight{\ifdim\Gin@nat@height>0.75\textheight0.75\textheight\else\Gin@nat@height\fi}
\makeatother
\setkeys{Gin}{width=\maxwidth,height=\maxheight,keepaspectratio}
\hypersetup{pdfauthor={OpenAI GPT-6.1 Sol, Codex, Ultra effort}}
\begin{document}
\section{Hilbert's Theorem 90 and Kummer
theory}\label{hilberts-theorem-90-and-kummer-theory}

\emph{Written by OpenAI GPT-6.1 Sol in Codex, Ultra effort, October
2026. Self-checked by the writing AI; no independent review is claimed.
Public domain (CC0).}

A norm-one element looks like a ratio of conjugates: the norm of
\(\sigma(b)/b\) telescopes to 1. Hilbert's Theorem 90 says that this
observation accounts for every norm-one element in a cyclic extension.
Its cocycle form works for any finite Galois group and supplies the
existence step in Kummer theory. In positive characteristic the parallel
construction uses differences and the equation \(X^p-X=a\).

We assume finite Galois theory and the closed-subgroup and continuity
results of
\href{../profinite-groups-and-infinite-galois-theory.html}{Profinite
groups and infinite Galois theory}. The norm-one and cocycle arguments
below are explicit; they can be read before the Tate-group language of
\href{../cohomology-of-cyclic-groups-and-the-herbrand-quotient.html}{Cohomology
of cyclic groups and the Herbrand quotient}. Basic references are
Milne's freely available notes and the Stacks project. All extensions
below lie in a chosen separable closure when a common ambient field is
needed.

\subsection{1. Detecting a nonzero averaging
operator}\label{detecting-a-nonzero-averaging-operator}

If \(L/K\) is finite Galois, averaging \(x\) over its conjugates gives a
trace. In characteristic dividing \([L:K]\), the trace of 1 is zero.
That does not mean that the trace map is zero. The following
independence argument will provide the nonzero averages we actually
need.

\textbf{Dedekind's independence lemma.} Distinct group homomorphisms
\(\chi_1,\ldots,\chi_m:S\to F^\times\), where \(F\) is a field, are
linearly independent over \(F\) as functions on \(S\). In particular,
distinct field embeddings into \(F\) are linearly independent.

\textbf{Proof.} Suppose a nonzero relation exists and choose one with as
few nonzero coefficients as possible. A one-term relation is impossible
because a character takes nonzero values. Relabel its terms so that all
its coefficients are nonzero. Choose \(h\in S\) with
\(\chi_1(h)\ne\chi_m(h)\). Evaluate the relation at \(hx\) and subtract
\(\chi_m(h)\) times the relation at \(x\). The last term disappears
while the first does not. This is a nonzero relation with fewer terms, a
contradiction. For embeddings, apply the result to their restrictions to
the multiplicative group of the domain field. \(\square\)

Thus a linear combination \(\sum_{g\in G}a_g g\) of distinct Galois
automorphisms is a nonzero operator whenever its coefficients are not
all zero. This is stronger than asserting that a particular value, such
as its value at 1, is nonzero.

\subsection{2. Ratios of conjugates}\label{ratios-of-conjugates}

\subsubsection{Theorem 3.1. Hilbert's Theorem
90}\label{theorem-3.1.-hilberts-theorem-90}

For a finite Galois extension \(L/K\), with group \(G\),

\[
H^1(G,L^\times)=0.
\]

If \(G=\langle\sigma\rangle\) is cyclic, then

\[
N_{L/K}(a)=1
\quad\Longleftrightarrow\quad
a=\frac{\sigma(b)}b\text{ for some }b\in L^\times.
\]

\textbf{Proof.} Let \(c:G\to L^\times\) be a multiplicative cocycle, so
\(c(gh)=c(g)g(c(h))\). Independence gives \(t\in L\) such that

\[
S=\sum_{h\in G}c(h)h(t)\ne0.
\]

For \(g\in G\), reindex the sum by \(gh\):

\[
g(S)=\sum_h g(c(h))gh(t)
=c(g)^{-1}\sum_h c(gh)gh(t)
=c(g)^{-1}S.
\]

Taking \(b=S^{-1}\) gives \(g(b)/b=c(g)\). Every cocycle is therefore a
coboundary.

For cyclic \(G\) of order \(n\), a norm-one element \(a\) defines the
cocycle

\[
c(\sigma^j)=\prod_{i=0}^{j-1}\sigma^i(a).
\]

The norm-one condition makes this definition compatible with
\(\sigma^n=1\). Apply the first part and evaluate at \(\sigma\).
Conversely, the norm of \(\sigma(b)/b\) is 1 by telescoping. \(\square\)

For example, in \(\mathbf Q(i)\), take \(\sigma(i)=-i\) and \(b=2-i\).
Then

\[
\frac{\sigma(b)}b=\frac{2+i}{2-i}=\frac{3+4i}{5}.
\]

Any two nonzero choices of \(b\) for the same ratio differ by a factor
in \(\mathbf Q^\times\): their quotient is fixed by \(\sigma\). The
theorem produces an element, not a preferred element.

\subsection{3. Differences of
conjugates}\label{differences-of-conjugates}

\subsubsection{Proposition 3.2. Additive Hilbert
90}\label{proposition-3.2.-additive-hilbert-90}

For a finite Galois extension \(L/K\),

\[
H^1(G,L)=0.
\]

If \(G=\langle\sigma\rangle\) is cyclic, then
\(\operatorname{Tr}_{L/K}(a)=0\) if and only if \(a=\sigma(b)-b\) for
some \(b\in L\).

\textbf{Proof.} Independence says that \(\sum_{g\in G}g\) is a nonzero
operator. Its image lies in \(K\). Choose \(t\) with
\(\operatorname{Tr}_{L/K}(t)=1\), by scaling a value with nonzero trace.

Given an additive cocycle \(c(gh)=c(g)+g c(h)\), put
\(S=\sum_h c(h)h(t)\). Reindexing gives

\[
g(S)=\sum_h(c(gh)-c(g))gh(t)=S-c(g).
\]

Thus \(b=-S\) satisfies \(g(b)-b=c(g)\). For a cyclic group of order
\(n\), a trace-zero element \(a\) defines the cocycle
\(c(\sigma^j)=\sum_{i=0}^{j-1}\sigma^i(a)\), with \(c(1)=0\). Trace zero
makes this formula periodic modulo \(n\); splitting a sum at \(j\)
verifies \(c(\sigma^{j+k})=c(\sigma^j)+\sigma^j c(\sigma^k)\), including
when the indices wrap around. The coboundary formula at \(\sigma\) gives
\(a=\sigma(b)-b\). Conversely that difference has trace zero by
telescoping. \(\square\)

This proof works even when the characteristic divides \(|G|\). Dividing
by \(|G|\) would fail in that case; choosing a trace-one element avoids
that division.

\textbf{Finite character lemma.} If a finite abelian group \(B\) is
killed by \(n\), and a field contains all \(n\) distinct \(n\)-th roots
of unity, then \(\operatorname{Hom}(B,\mu_n)\) has \(|B|\) elements and
its characters separate elements of \(B\).

\textbf{Proof.} Write \(B\) additively. Every finite subgroup of a
field's multiplicative group is cyclic, by the exponent and
polynomial-root argument of lesson 1, section 0. In particular \(\mu_c\)
is cyclic of order \(c\) for every \(c\mid n\). Suppose a character
\(\chi\) is defined on a subgroup \(A\subset B\), and adjoin \(x\in B\).
Let \(c\) be the order of \(x\), and \(d\) its order modulo \(A\); then
\(d\mid c\) and \(dx\in A\). We have \(\chi(dx)^{c/d}=1\). In the cyclic
group \(\mu_c\), the image of the \(d\)-th power map is exactly
\(\mu_{c/d}\), so choose \(z\in\mu_c\) with \(z^d=\chi(dx)\). The
formula \[
\chi'(a+kx)=\chi(a)z^k
\] defines an extension to \(A+\langle x\rangle\): changing a
representation changes \(k\) by a multiple of \(d\), and the
accompanying factor from \(A\) cancels by the equation for \(z\). There
are exactly \(d\) choices for \(z\) in \(\mu_n\), since \(d\mid n\); all
have \(z^c=(z^d)^{c/d}=1\), so each gives an extension. Starting at the
trivial subgroup and adjoining generators therefore multiplies the
character count by the same indices as the group order. The count is
\(|B|\). To separate a nonzero \(x\), first assign to \(x\) a primitive
root of its order on \(\langle x\rangle\), then extend this character
through the remaining generators. Its value at \(x\) stays nontrivial.
\(\square\)

\subsection{4. Radicals as characters}\label{radicals-as-characters}

Let \(n\geq1\) be prime to \(\operatorname{char}K\), and assume
\(\mu_n\subset K\). Write
\(G_K=\operatorname{Gal}(K^{\mathrm{sep}}/K)\). For \(a\in K^\times\),
choose \(\alpha\) with \(\alpha^n=a\), and define

\[
\chi_a(g)=\frac{g(\alpha)}{\alpha}\in\mu_n.
\]

All \(n\)-th roots of \(a\) differ by an element of \(\mu_n\subset K\),
so this character does not depend on the chosen root. It is a
homomorphism because \(G_K\) acts trivially on \(\mu_n\), and it is
continuous because \(\alpha\) lies in a finite extension. Multiplying
\(a\) by an \(n\)-th power in \(K\) does not change it.

The resulting map

\[
K^\times/K^{\times n}\longrightarrow
\operatorname{Hom}_{\mathrm{cont}}(G_K,\mu_n)
\]

is an isomorphism. Injectivity holds because a root fixed by all of
\(G_K\) belongs to \(K\). For surjectivity, a continuous character
factors through the Galois group of some finite Galois \(E/K\): its open
kernel contains an open normal subgroup. Regard the character as an
\(E^\times\)-valued cocycle. Hilbert 90 gives \(\alpha\in E^\times\)
with \(g(\alpha)/\alpha=\chi(g)\). Then \(\alpha^n\) is fixed by the
finite group and belongs to \(K^\times\), giving the desired preimage.

\subsubsection{Theorem 3.3. Kummer
correspondence}\label{theorem-3.3.-kummer-correspondence}

Under these hypotheses, the assignments

\[
\Delta\longmapsto K(\Delta^{1/n}),
\qquad
L\longmapsto\{a\in K^\times:a\text{ has an }n\text{-th root in }L\}
\]

are inverse inclusion-preserving bijections between subgroups
\(K^{\times n}\subset\Delta\subset K^\times\) and abelian Galois
extensions \(L/K\) of exponent dividing \(n\). They include infinite
extensions. The pairing

\[
\langle g,a\rangle=\frac{g(\alpha)}\alpha,
\qquad \alpha^n=a,
\]

identifies

\[
\operatorname{Gal}(K(\Delta^{1/n})/K)
\simeq\operatorname{Hom}(\Delta/K^{\times n},\mu_n).
\]

Give \(\Delta/K^{\times n}\) the discrete topology and the Hom group the
topology inherited from the product \(\mu_n^{\Delta/K^{\times n}}\).
Finite subgroups correspond to finite extensions, with

\[
[K(\Delta^{1/n}):K]=|\Delta/K^{\times n}|.
\]

\textbf{Proof.} Start with a finite subgroup \(D=\Delta/K^{\times n}\).
Choose finitely many representatives generating it. Their root fields
are Galois because all the required roots of unity are in \(K\); their
compositum \(L\) is finite Galois. Its Galois action injects into
\(\operatorname{Hom}(D,\mu_n)\), since the roots generate \(L\). In
particular its group \(G\) is abelian of exponent dividing \(n\).

The pairing also injects \(D\) into \(\operatorname{Hom}(G,\mu_n)\): if
every element of \(G\) fixes a chosen root, that root is in \(K\), so
the represented class is zero. For every finite abelian group \(B\)
killed by \(n\),

\[
|\operatorname{Hom}(B,\mu_n)|=|B|.
\]

This is the finite character lemma proved above. The two injections
consequently force \(|G|=|D|\), and both pairing maps are isomorphisms.
This proves the degree formula and perfectness.

For a finite abelian extension \(L/K\) of exponent dividing \(n\), each
character \(G\to\mu_n\) is an \(L^\times\)-valued cocycle. Hilbert 90
realizes it by a radical \(\alpha\in L\), with \(\alpha^n\in K\). The
finite character lemma says these characters separate elements. Hence
the common stabilizer of these radicals is trivial, and they generate
\(L\). Moreover, \(D_L=(L^{\times n}\cap K^\times)/K^{\times n}\) maps
bijectively to \(\operatorname{Hom}(G,\mu_n)\) by this same argument.
For an extension originally constructed from \(D\), the subgroup
\(D\subset D_L\) already has order \(|G|\), so it equals \(D_L\). The
two assignments are inverse in the finite case.

For arbitrary \(D\), take the union of its finitely generated subgroups.
Each is finite, since its generators have order dividing \(n\). Their
root fields have union \(K(\Delta^{1/n})\), and the inverse-limit
description of its Galois group is exactly
\(\operatorname{Hom}(D,\mu_n)\). Conversely, an infinite abelian
extension of exponent dividing \(n\) is the union of its finite Galois
subextensions. The finite correspondence applies to each. A root in the
union lies in one finite subextension, so the inverse assignment
introduces no additional classes. This proves both the infinite
correspondence and the topology assertion. \(\square\)

For finite \(D\), perfectness means that each side is the full dual of
the other. For infinite \(D\), the second dual consists of the
\textbf{continuous} characters of its compact Galois group: a continuous
character factors through a finite stage. Omitting this continuity would
give a larger and incorrect dual.

If \(L/K\) is cyclic of degree \(n\), choose a faithful character
\(G\to\mu_n\) and apply Hilbert 90. Its radical has trivial stabilizer
and generates \(L\), proving \(L=K(a^{1/n})\). Conversely,
\(K(a^{1/n})/K\) is cyclic of degree the order of the class of \(a\) in
\(K^\times/K^{\times n}\); that order can be a proper divisor of \(n\).

For example, the three independent square classes 2, 3 and 5 give

\[
\mathbf Q(\sqrt2,\sqrt3,\sqrt5)/\mathbf Q,
\qquad G\simeq(\mathbf Z/2\mathbf Z)^3.
\]

Their independence follows from the valuations at 2, 3 and 5. The
pairing records the three independent sign changes.

\subsection{5. The additive version in characteristic
p}\label{the-additive-version-in-characteristic-p}

Assume now \(\operatorname{char}K=p>0\), and put \(\wp(x)=x^p-x\). This
is an additive map whose kernel is \(\mathbf F_p\). Every polynomial
\(X^p-X-a\) is separable, since its derivative is \(-1\).

For a root \(\alpha\), the function \(g\mapsto g(\alpha)-\alpha\) takes
values in \(\mathbf F_p\) and is a continuous homomorphism. Changing the
root adds a constant in \(\mathbf F_p\), and changing \(a\) by
\(\wp(b)\), for \(b\in K\), replaces a root by \(\alpha+b\). Thus this
construction depends only on \(a\bmod\wp(K)\).

\subsubsection{Proposition 3.4. Artin--Schreier
correspondence}\label{proposition-3.4.-artinschreier-correspondence}

There is an isomorphism of \(\mathbf F_p\)-vector spaces

\[
K/\wp(K)\simeq\operatorname{Hom}_{\mathrm{cont}}(G_K,\mathbf F_p).
\]

Subspaces \(D\subset K/\wp(K)\) correspond, preserving inclusion, to
abelian Galois extensions of exponent dividing \(p\), by adjoining roots
of \(X^p-X-a\) for the classes \(a\in D\). Their Galois groups are

\[
\operatorname{Hom}_{\mathbf F_p}(D,\mathbf F_p),
\]

with the product topology. A finite-dimensional \(D\) gives degree
\(p^{\dim D}\).

\textbf{Proof.} A class maps to zero exactly when its root is fixed by
\(G_K\), hence lies in \(K\); that is exactly \(a\in\wp(K)\). A
continuous homomorphism to \(\mathbf F_p\) factors through some finite
Galois \(E/K\). Additive Hilbert 90 realizes it as \(g(\alpha)-\alpha\),
with \(\alpha\in E\). Then \(\wp(\alpha)\) is fixed and lies in \(K\),
proving surjectivity.

For a finite-dimensional subspace, its root field is finite Galois, and
its action injects into the vector-space dual of \(D\). The translation
pairing injects \(D\) into the dual of the Galois group, since a root
fixed by that group lies in \(K\). Equal dimensions, or the two
resulting cardinality inequalities, make both maps isomorphisms.
Conversely, additive Hilbert 90 realizes all characters of a finite
elementary abelian \(p\)-group by roots; they separate the group and
therefore generate its field. This proves the finite correspondence and
its inverse. Taking unions of finite-dimensional subspaces and inverse
limits of their Galois groups proves the general case, exactly as in
Theorem 3.3. \(\square\)

For \(K=\mathbf F_p(t)\), the class of \(t^{-1}\) is nonzero in
\(K/\wp(K)\). A pole of order \(m>0\) in \(b\) becomes a pole of order
\(pm\) in \(b^p-b\); if \(b\) has no pole at \(t=0\), neither does
\(b^p-b\). Neither possibility gives the pole of order 1 of \(t^{-1}\).
Hence \(X^p-X=t^{-1}\) defines a cyclic extension of degree \(p\).

\subsection{Exercises}\label{exercises}

\begin{enumerate}
\def\labelenumi{\arabic{enumi}.}
\tightlist
\item
  \textbf{Easy.} In \(\mathbf Q(i)\), find every \(b\ne0\) satisfying
  \(\sigma(b)/b=(3+4i)/5\), where \(\sigma(i)=-i\).
\item
  \textbf{Medium.} Describe all quadratic and biquadratic extensions of
  \(\mathbf Q\) in terms of square classes. Apply the description to the
  two classes \(-2\) and 7.
\item
  \textbf{Medium.} For a finite subgroup
  \(D\subset K^\times/K^{\times n}\), prove that the Kummer pairing of
  Theorem 3.3 is perfect without assuming a degree formula for radical
  extensions.
\item
  \textbf{Hard.} Prove additive Hilbert 90 in arbitrary characteristic
  using a trace-one element. Use it to derive the Artin--Schreier
  character isomorphism and the correspondence for infinite elementary
  abelian extensions.
\end{enumerate}

\subsection{Solutions}\label{solutions}

\begin{enumerate}
\def\labelenumi{\arabic{enumi}.}
\tightlist
\item
  The element \(b_0=2-i\) satisfies the equation by direct
  multiplication. If \(b\) is another solution, then
  \(\sigma(b/b_0)=b/b_0\), so \(b/b_0\in\mathbf Q^\times\). All
  solutions are therefore \(b=r(2-i)\), with \(r\in\mathbf Q^\times\).
\item
  A quadratic extension is determined by a one-dimensional subspace of
  \(\mathbf Q^\times/\mathbf Q^{\times2}\), and has the form
  \(\mathbf Q(\sqrt d)\), where \(d\) is a nontrivial square class. It
  has a unique squarefree integer representative. A biquadratic
  extension is determined by a two-dimensional subspace, generated by
  two independent classes \(a,b\); its field is
  \(\mathbf Q(\sqrt a,\sqrt b)\), and its three quadratic subfields
  correspond to \(a,b,ab\). Different bases of the same subspace give
  the same extension. The classes \(-2\) and 7 are independent by their
  valuations at 2 and 7. The resulting field has degree 4 and quadratic
  subfields \(\mathbf Q(\sqrt{-2})\), \(\mathbf Q(\sqrt7)\) and
  \(\mathbf Q(\sqrt{-14})\).
\item
  Let \(L\) be the root field and \(G=\operatorname{Gal}(L/K)\). The map
  \(G\to\operatorname{Hom}(D,\mu_n)\) is injective because the roots
  generate \(L\). This also proves that \(G\) is abelian of exponent
  dividing \(n\). The map in the other direction is injective because a
  root fixed by all of \(G\) lies in \(K\). The finite character lemma
  above proves that each of these dual groups has the same order as its
  group. Thus \(|G|\leq|D|\leq|G|\), and both injections are surjective.
  This proves perfectness and simultaneously establishes the degree
  formula.
\item
  Independence of the distinct embeddings makes the trace a nonzero
  \(K\)-linear map to \(K\), so choose \(t\) of trace 1. For an additive
  cocycle, \(S=\sum_h c(h)h(t)\) satisfies \(gS=S-c(g)\), and \(b=-S\)
  gives \(c(g)=gb-b\). No division by \(|G|\) occurs. Given a continuous
  character \(G_K\to\mathbf F_p\), factor it through finite Galois
  \(E/K\) and use this construction to obtain \(\alpha\) with
  \(g\alpha-\alpha=\chi(g)\). Then \(\alpha^p-\alpha\in K\). Its class
  maps to the character; the kernel consists exactly of \(\wp(K)\). The
  finite-dimensional translation pairing is perfect by the two
  injections and equality of dual dimensions. Any subspace is a union of
  finite-dimensional ones; their root fields have the corresponding
  union, and their Galois groups have the inverse limit. Conversely,
  every element of an infinite elementary abelian extension lies in a
  finite Galois subextension. Applying the finite inverse assignments
  there proves the infinite correspondence with no finiteness assumption
  on \(K/\wp(K)\).
\end{enumerate}

\subsection{Editable edition}\label{editable-edition}

The reading edition provides the complete LaTeX source of this lesson,
the cumulative course LaTeX and the editable source ZIP. The archive
contains all twenty-four Markdown lessons, complete LaTeX bodies,
original diagrams, metadata and reproduction instructions.

\subsection{References}\label{references}

Sections 1--3 prove independence and both forms of Hilbert 90 in every
characteristic. The finite character lemma proves dual size and
separation before the full finite and infinite Kummer and
Artin--Schreier correspondences.

\begin{itemize}
\tightlist
\item
  \href{https://www.jmilne.org/math/CourseNotes/FT.pdf}{J. S. Milne,
  Fields and Galois Theory, version 5.10}.
\item
  \href{https://stacks.math.columbia.edu/}{The Stacks project,
  independence, Kummer and Artin--Schreier sections}.
\end{itemize}

The Stacks comparisons are Tags
\href{https://stacks.math.columbia.edu/tag/0CKK}{0CKK},
\href{https://stacks.math.columbia.edu/tag/09I6}{09I6} and
\href{https://stacks.math.columbia.edu/tag/09I7}{09I7}. Related
independently written programme material belongs to the AI Integrated
Stacks Project; it has not been reviewed by the official Stacks
maintainers.

The \href{../free-proofs.html}{proof guide} gives the lesson sequence
and the exact prerequisite record. External references accompany the
written arguments.

\end{document}
